\documentclass[10pt,a4paper]{amsart}
\usepackage[margin=19mm]{geometry}
\usepackage{amsmath,amssymb,mathtools}
\usepackage[scr=rsfs,cal=euler]{mathalfa}
\usepackage{xfrac}
\usepackage[unicode,hidelinks,bookmarks=false]{hyperref}
\newcommand{\ie}{i.e.\ }
\newcommand{\cf}{cf.\ }
\newcommand{\resp}{respectively\ }
\newcommand{\Subsection}{\subsection}
\providecommand{\nobreakdash}{\nobreak\hbox{-}}
\setlength{\emergencystretch}{2em}
\multlinegap0pt
\usepackage{amssymb}
\usepackage{dsfont}
\newtheorem{theorem}{Theorem}
\newtheorem{lemma}{Lemma}
\newcommand{\graph}{\mathcal{G}}
\newcommand{\plim}[0]{p\mbox{-}\lim}
\newcommand{\prob}[1]{\mathbb{P}\left(#1\right)}
\newcommand{\tti}{\rightarrow\infty}
\title{Entropy and Distributed Source Coding of Connected Soft Random Geometric Graphs}
\author{Oliver Baker, Carl P. Dettmann}
\date{Source: arXiv:2605.04703v2 (CC BY 4.0)}
\begin{document}
\maketitle

\noindent\emph{Source and licence.} Entropy and Distributed Source Coding of Connected Soft Random Geometric Graphs. Version arXiv:2605.04703v2 (CC BY 4.0), distributed by arXiv under the Creative Commons Attribution 4.0 International licence, http://creativecommons.org/licenses/by/4.0/, as recorded in the arXiv licence metadata for this exact version and shown on the abstract page https://arxiv.org/abs/2605.04703v2. Full licence notice: \texttt{licenses/arxiv-2605.04703-CC-BY-4.0.txt}.\par\smallskip

\noindent\textit{Editorial scope.} Counted unit: the article's lemma thm:nbhd\_reduction (source0.tex lines 357--368). The article's complete original proof of it is delivered with it (source0.tex lines 370--383, one \texttt{proof} environment of 14 lines). No mathematics is written, repaired or re-derived: the statement and the proof are the article's own lines, in the article's own notation, and no retained line is substituted or re-flowed.  Retained uncounted as the source-local context the proof needs: lines 224--247.  Nothing was omitted from any retained span, so the package carries no omission record.  No retained line contains a citation, so the wrapper carries no bibliography and no bibliography entry.  No retained line contains a figure environment, an image-inclusion command or drawing code, so no image is delivered and the package declares no asset dependency.  Editorial formatting only, in addition: an AMS \texttt{amsart} preamble that restates the article's own declarations and macros used by the retained lines, namely the environments \texttt{theorem}, \texttt{lemma} on separate counters, together with the standard packages those lines need.  The retained environments print in this document as Theorem 1, Lemma 1 in order of appearance, whereas the article prints the same items with the numbering of its own full document.  The complete native source of the article as distributed in the arXiv e-print for this exact version is bundled unchanged as \texttt{sources/199887/source0.tex}.
\begin{theorem}[Asymptotic Equipartition Property]
    \label{thm:aep}
    If $G_n$ is a random SRGG sampled from $\graph_n$,
    \begin{equation}
        \plim_{n\tti} \frac{1}{\binom{n}{2}s(n)^d} \log \frac{1}{\prob{G_n}} = h^*
    \end{equation} 
    In other words,
    \begin{equation}
        \overline{H}(\mathbf{G}) = \underline{H}(\mathbf{G}) = H(\mathbf{G}) = h^*
    \end{equation}
    This implies the existence of sets defined for all $\epsilon > 0$,
    \begin{equation}
        \mathcal{T}_n^\epsilon := \left\{G_n : \left|\log \frac{1}{\prob{G_n}} - H(\graph_n)\right| \leq \binom{n}{2}s(n)^d \epsilon\right\} 
    \end{equation}
    such that
    \begin{enumerate}
        \item $\prob{G_n \in \mathcal{T}_n^\epsilon} \to 1$ as $n\tti$
        \item For each $G_n \in \mathcal{T}_n^\epsilon$, $2^{-\binom{n}{2}(s(n)^dh^* +\epsilon)} \leq \prob{G_n} \leq 2^{-\binom{n}{2}s(n)^d(h^*- \epsilon)}$
        \item For every $\delta > 0$, there exists an $n^*$ so that for all $n > n^*$, we have
        \begin{equation}
            (1-\delta)2^{\binom{n}{2}s(n)^d(h^*-\epsilon)} \leq |\mathcal{T}_n^\epsilon| \leq 2^{\binom{n}{2}s(n)^d(h^*+\epsilon)}
        \end{equation}
    \end{enumerate}
\end{theorem}

\begin{lemma}
    \label{thm:nbhd_reduction}
    Let $S_n \subseteq V_n$, and $S_n^c := V_n \setminus S_n$ then
    \begin{equation}
        H(N_{S_n}|N_{S_n^c}) = H(\graph_n[S_n])
    \end{equation}
    where $\graph_n[S_n]$ is the ensemble of subgraphs $G_n[S_n]$ of $G_n$ restricted to node set $S$. This implies that if $\{S_n\}_{n=1}^\infty$ is a sequence of subsets of $V_n$ indexed by $n$, such that $|S_n| \to \infty$ as $n\tti$, then
    \begin{equation}
        \label{eq:nbhd_aep}
        \plim_{n\tti} \frac{1}{\binom{|S_n|}{2}s(n)^d}\log \frac{1}{\prob{N_{S_n}|N_{S_n^c}}} = h^* 
    \end{equation}
\end{lemma}

\begin{proof}
    For the first part, fix $n$, and let $k < n$. We will show the first part for ${S} = \{1,..,k\}$,  which generalises to any set ${S}$ by re-labelling the graph. Writing out $N_{{S}}$ in full, we have, recalling that $X_{ij}$ is identified with $X_{ji}$,
    \begin{equation}
        N_{{S}} = \{X_{ij}\}_{\substack{i \in {S}, 1 \leq j \leq n\\
                                    i \neq j}}, N_{{S}^c} = \{X_{ij}\}_{\substack{i \notin {S}, 1 \leq j \leq n\\
                                    i \neq j}}
    \end{equation}
    Then, the entropy $H(N_{S}|N_{S^c
    })$ is the entropy of $A_{{S}} := N_{S} \setminus (N_{S} \cap N_{{S}^c})$ conditioned on $\widetilde{A}_{S} := N_{{S}^c} \setminus (N_{{S}} \cap N_{{S}^c})$, since the edges in the intersection do not contribute to the entropy. The key observation is that if $X_{ij} \in A_{S}$, then there does not exist an edge $X_{uv}$ in $\widetilde{A}_{S}$ with $u = i$ or $v = j$. Therefore, every edge variable in $A_{S}$ is independent of all those in $\widetilde{A}_{S}$. Then,
    \begin{equation*}
        H(N_S|N_{S^c}) = H(A_{S}|\widetilde{A}_{S}) = H(\{X_{ij}\}_{1\leq i<j \leq k}) = H(\graph_n[S])
    \end{equation*}
    as required. For the second part of the theorem, let $\{S_n\}$ be a sequence of sets with size proportional to $n$. Since each edge is added to the graph independently, the subgraph $G_n[S]$ is marginally distributed as a SRGG on $|S|$ nodes with sparsity $s(n)$ inherited from the parent graph $G_n$. Then, we can apply Theorem \ref{thm:aep} to get the result.
\end{proof}

\end{document}
